\documentclass{article}

% Recommended packages
\usepackage{microtype}
\usepackage{graphicx}
\usepackage{subcaption}
\usepackage{booktabs}
\usepackage{hyperref}
\newcommand{\theHalgorithm}{\arabic{algorithm}}

\usepackage{icml2026}

\usepackage{amsmath}
\usepackage{amssymb}
\usepackage{mathtools}
\usepackage{amsthm}
\usepackage[capitalize,noabbrev]{cleveref}

% Theorems
\theoremstyle{plain}
\newtheorem{theorem}{Theorem}[section]
\newtheorem{proposition}[theorem]{Proposition}
\newtheorem{lemma}[theorem]{Lemma}
\newtheorem{corollary}[theorem]{Corollary}
\theoremstyle{definition}
\newtheorem{definition}[theorem]{Definition}
\newtheorem{assumption}[theorem]{Assumption}
\theoremstyle{remark}
\newtheorem{remark}[theorem]{Remark}

\icmltitlerunning{Monte Carlo Variational Task-Information Merging}

\begin{document}

\twocolumn[
  \icmltitle{MC-VTI: Monte Carlo Variational Task-Information Merging for \\
              Stabilized Test-Time Model Merging}

  \icmlsetsymbol{equal}{*}

  \begin{icmlauthorlist}
    \icmlauthor{Anonymous Authors}{equal}
  \end{icmlauthorlist}

  \icmlaffiliation{equal}{Under Review for ICML 2026}
  \icmlcorrespondingauthor{Anonymous Authors}{authors@anon.org}

  \icmlkeywords{Model Merging, Test-Time Adaptation, Bayesian Neural Networks, Uncertainty Estimation}

  \vskip 0.3in
]

\printAffiliationsAndNotice{}

\begin{abstract}
  Model merging has emerged as a highly cost-effective paradigm to integrate specialized capabilities of independently fine-tuned expert models into a single multi-task system without requiring prohibitive joint retraining. To handle non-stationary and dynamically shifting environments, recent works adapt merging weights on unlabeled test streams via test-time adaptation (TTA). However, unsupervised test-time adaptation is notoriously susceptible to overfitting, parameter drift, and decision-boundary collapse in classification heads. Existing methods either freeze classification heads completely (violating adaptability under domain shifts) or rely on pre-computing Fisher Information Matrix (FIM) priors on the original training datasets (violating privacy and data availability). In this work, we propose \textbf{MC-VTI} (Monte Carlo Variational Task-Information Merging), a fully self-supervised, teacher-free, and training-data-free Bayesian TTA framework for model merging. MC-VTI utilizes test-time Monte Carlo Dropout to approximate parameter posteriors. It optimizes expected prediction entropy over multiple MC dropout passes to locate flat, robust loss basins, and incorporates an online translation-consistency regularizer. Crucially, MC-VTI computes the logit variance across MC passes on-the-fly and applies a dynamic quadratic penalty to the classification heads, protecting parameters of highly uncertain classes from catastrophic drift. Rigorous evaluations on a multi-task MNIST-FashionMNIST-KMNIST ResNet-18 benchmark demonstrate that MC-VTI completely prevents decision boundary collapse, outperforming frozen-head baselines and performing competitively with methods requiring offline training-data priors.
\end{abstract}

\section{Introduction}
\label{sec:intro}
The pre-training and fine-tuning paradigm has driven deep learning to massive success across diverse applications in computer vision, natural language processing, and multimodal representation learning. Typically, a large foundational model is pre-trained on a massive, general-world dataset and subsequently fine-tuned on specific downstream tasks to specialize its capabilities. However, fine-tuning large base models independently on multiple downstream tasks yields many separate experts, scaling parameter storage, hosting, and edge-deployment costs linearly with the number of tasks. Joint multi-task training (e.g., training a single model on all datasets simultaneously) bypasses this storage explosion, but it is often computationally prohibitive, requires simultaneous access to all private datasets, and introduces severe optimization difficulties such as gradient conflicts, negative transfer, and Pareto-frontier compromises \cite{yadav2023ties,chen2024deep}.

To circumvent these limitations, model merging (or deep model fusion) has emerged as an exceptionally attractive, training-free alternative \cite{wortsman2022model,ilharco2023editing,yang2024bridging,amine2019merging,wang2024dynamic}. Model merging directly combines independently fine-tuned weights at the parameter level, integrating multi-task specialized capabilities without requiring prohibitive joint retraining. Traditional methods such as Task Arithmetic \cite{ilharco2023editing} and TIES-Merging \cite{yadav2023ties} perform simple linear operations in Euclidean space to merge task-specific vectors. While highly efficient, these Euclidean averaging techniques suffer from severe parameter interference, where conflicting updates from different experts cancel each other out, degrading performance. Furthermore, static merging coefficients are incapable of adapting to distribution shifts or non-stationary streams at test time, limiting their utility in real-world environments.

To address parameter interference and provide dynamically adaptive routing, test-time adaptive merging methods optimize merging coefficients on unlabeled test data on the fly \cite{sata2026convex,s2c2026teacher,ewc2026elastic}. However, real-world deployment exposes systems to non-stationary, shifting, or severely biased test streams where the underlying task distribution changes dynamically and abruptly over time (e.g., switching from digital digits to fashion accessories). Standard unconstrained test-time adaptation (TTA) methods applied to model merging can alleviate these limitations by dynamically adjusting merging coefficients, but they suffer from severe representation collapse and decision-boundary drift in highly sensitive classification heads under severe stream distribution shifts. Specifically, because the incoming test streams are unlabeled, standard self-supervised objectives like entropy minimization can overfit to noisy local batches, leading to catastrophic drift in the classification parameters.

To protect the classification heads and stabilize test-time model merging, existing methods employ distinct strategies. For instance, S2C-Merge \cite{s2c2026teacher} completely freezes the classification heads and restricts adaptation strictly to a low-dimensional merging coefficient space. While this prevents decision boundary collapse, it fundamentally prevents the classification heads from adapting to out-of-distribution (OOD) test-time shifts and local covariates, restricting overall model adaptability. On the other hand, EWC-TTA \cite{ewc2026elastic} pre-computes diagonal Fisher Information Matrix (FIM) priors on clean training or validation datasets of the experts, using an Elastic Weight Consolidation (EWC) style quadratic penalty to protect task-critical parameter structures. However, in many real-world scenarios, the original training datasets are unavailable due to licensing restrictions, strict user privacy regulations (e.g., GDPR), or local edge-storage constraints, making EWC-TTA completely inapplicable in data-constrained or private environments.

To resolve these critical limitations, we propose \textbf{MC-VTI} (Monte Carlo Variational Task-Information Merging), a mathematically rigorous, fully self-supervised, teacher-free, and training-data-free test-time adaptation framework for model merging. MC-VTI leverages test-time Monte Carlo (MC) Dropout to construct a variational approximation of the model's parameter posterior distribution. By taking the expectation of the prediction entropy and translation consistency across multiple MC passes, MC-VTI optimizes a robust variational objective that directs updates toward flat, noise-resilient loss basins, mitigating overfitting to noisy online batches. 

Crucially, rather than relying on pre-computed training-set priors or freezing parameters, MC-VTI estimates model uncertainty on-the-fly by computing the MC logit variance for each class. It then applies a dynamic, uncertainty-weighted quadratic regularization penalty directly to the classification head parameters. For classes where the logit variance across MC passes is high (indicating high epistemic uncertainty and unstable boundaries), the penalty is strong, locking the decision boundaries and preventing representation collapse. For classes where the variance is low (indicating high certainty and stable representation structures), the penalty is relaxed, allowing the classification heads to adapt gently to OOD test data and local covariates. 

In summary, our key contributions are:
\begin{itemize}
  \item We introduce MC-VTI, a novel Bayesian-guided test-time adaptation framework for model merging that is completely teacher-free and training-data-free, resolving both adaptability and privacy challenges.
  \item We formulate an expected prediction entropy and translation-consistency objective under test-time MC Dropout to locate flatter, more robust parameter states that generalize well.
  \item We propose a dynamic, uncertainty-weighted regularization mechanism based on MC logit variance to stabilize classification heads on-the-fly, establishing a formal equivalence to task-specific Gaussian priors.
  \item We demonstrate that MC-VTI completely prevents decision boundary collapse, outperforming static and frozen-classifier baselines on multi-task vision benchmarks under sequential and alternating test streams, while providing honest scientific insights into the fundamental limits of gradient-based TTA under high-frequency switching.
\end{itemize}

\section{Related Work}
\label{sec:related}
\textbf{Model Merging:} Model merging aims to combine multiple independent, specialized models fine-tuned from a shared base model into a single multi-task model without further training, thereby bypassing joint multi-task training overheads \cite{wortsman2022model,ilharco2023editing,yadav2023ties,yu2024language}. Early approaches like Model Soups \cite{wortsman2022model} and Task Arithmetic \cite{ilharco2023editing} average parameters or scale and add task vectors. TIES-Merging \cite{yadav2023ties} addresses parameter interference by keeping only the top-$k$ most significant updates and resolving sign conflicts. DARE \cite{yu2024language} prunes small parameters randomly to maintain performance. Other recent advancements include combinatorial merging \cite{slimi2024combinatorial}, decentralized merges \cite{saadati2024decentralized}, and modular merges \cite{ladwig2023modular}. However, these static methods suffer from representation interference under severe domain shifts and are incapable of adapting dynamically to non-stationary test streams during deployment.

\textbf{Test-Time Adaptation (TTA):} TTA aims to adapt a pre-trained model to a target distribution on-the-fly during inference using unlabeled test streams. Tent \cite{wang2021tent} optimizes batch normalization parameters via prediction entropy minimization. SHOT \cite{liang2020source} freezes classifier heads and optimizes representations using information maximization. CoTTA \cite{wang2022continual} and RoTTA \cite{dobler2022robustness} handle non-stationary, continual streams using memory banks, self-distillation, and strong augmentations. Fully test-time adaptation \cite{wang2021fully,niu2023towards} and source-free domain adaptation (SFDA) \cite{yang2021exploiting,yang2022attracting,yang2021generalized,liu2021domain} optimize representations without accessing the original source datasets, which relates closely to our privacy and data availability constraints. Our method extends these paradigms by applying test-time Monte Carlo Dropout \cite{gal2015bayesian,hermosilla2021monte} to robustify adaptation and estimate uncertainty on-the-fly.

\textbf{Test-Time Adaptive Model Merging:} Integrating TTA into model merging allows the merging coefficients to adapt dynamically to the test stream. SATA-SBF \& SATA-RGP \cite{sata2026convex} utilize convex geometric constraints on merging weights and relative geometry preservation. S2C-Merge \cite{s2c2026teacher} proposes teacher-free TTA by optimizing layer-wise merging coefficients using self-supervised entropy and consistency. However, S2C-Merge freezes classification heads completely to prevent decision-boundary collapse, which limits OOD generalization. EWC-TTA \cite{ewc2026elastic} allows adapting both coefficients and heads by applying a quadratic penalty based on offline diagonal Fisher Information. In contrast to these, our MC-VTI is fully self-supervised, teacher-free, and training-data-free, estimating parameter uncertainty online to guide head stabilization on-the-fly, establishing a highly practical and secure foundation for deployment. We also ground our framework in Bayesian neural networks \cite{zhou2025bayesian} and variational inference \cite{gelberg2024variational,hubin2023variational} to justify our uncertainty estimation.

\section{Preliminaries \& Problem Setup}
\label{sec:prelims}
Let $f_{\theta_{pre}}$ be a pre-trained base model backbone (e.g., ResNet-18) parameterized by $\theta_{pre}$. We consider a set of $K$ distinct tasks $\{T_1, \dots, T_K\}$. For each task $k$, an expert model is independently fine-tuned, yielding backbone parameters $\theta_k$ and a task-specific classification head $h_{\phi_k}$ parameterized by $\phi_k$. The task-specific vector (or update delta) representing the acquired specialized knowledge is defined as:
\begin{equation}
  \tau_k = \theta_k - \theta_{pre}
\end{equation}
A merged model combines these updates using a set of merging coefficients $\lambda = [\lambda_1, \dots, \lambda_K]^T \in \mathbb{R}^K$. The merged backbone parameters $\theta(\lambda)$ are reconstructed on-the-fly:
\begin{equation}
  \theta(\lambda) = \theta_{pre} + \sum_{k=1}^K \lambda_k \tau_k
\end{equation}
During online test-time adaptation, the model receives a continuous stream of unlabeled batches $X_t = \{x_{t,1}, \dots, x_{t,B}\}$. At each step $t$, the task origin $k$ of the batch is unknown, but the active task head corresponding to the stream is evaluated. The goal is to optimize $\lambda$ and $\phi_k$ on-the-fly to minimize prediction error without accessing the original training data or keeping unmerged expert models in memory.

This problem setup introduces three critical challenges: (1) **Parameter Interference**: combining expert weights in parameter space can lead to representation degradation and negative transfer across tasks. (2) **Catastrophic Forgetting \& Parameter Drift**: adapting parameters using unlabeled data in a sequential or non-stationary stream can cause classification heads to drift, leading to decision boundary collapse. (3) **Data Constraints \& Privacy**: in many real-world environments, the original expert training datasets are private or too large to store on edge devices, preventing the pre-computation of static priors. MC-VTI addresses all three of these challenges by estimating model uncertainty online and regularizing the parameters dynamically.

\section{Methodology: MC-VTI}
\label{sec:method}
\begin{algorithm}[tb]
   \caption{MC-VTI: Monte Carlo Variational Task-Information Merging}
   \label{alg:mcvti}
\begin{algorithmic}[1] \footnotesize
   \STATE {\bfseries Input:} Base model $\theta_{pre}$, Expert backbones $\{\theta_k\}_{k=1}^K$, Expert heads $\{\phi_k\}_{k=1}^K$, Test stream $\{X_t\}_{t=1}^T$, learning rates $\eta_\lambda, \eta_\phi$, regularization strength $\gamma$, MC passes $M$
   \STATE {\bfseries Initialize:} Merging coefficients $\lambda \leftarrow [1/K, \dots, 1/K]^T$
   \STATE {\bfseries Extract:} Expert updates $\tau_k \leftarrow \theta_k - \theta_{pre}$ and save initial heads $\phi^{init}_k \leftarrow \text{copy}(\phi_k)$
   \FOR{each test step $t = 1$ to $T$}
       \STATE Receive unlabeled batch $X_t$ and active task index $task\_idx$
       \STATE Get active head parameters $\phi \leftarrow \phi_{task\_idx}$ and its initial state $\phi^{init} \leftarrow \phi^{init}_{task\_idx}$
       \STATE Reconstruct backbone parameters: $\theta(\lambda) \leftarrow \theta_{pre} + \sum_{k=1}^K \lambda_k \tau_k$
       \STATE Evaluate predictions: $z \leftarrow \text{functional\_call}(f_{\theta_{pre}}, \theta(\lambda), X_t)$ and $logits \leftarrow \phi(z)$
       \STATE // {\itshape Test-time variational adaptation step}
       \FOR{$m = 1$ to $M$}
           \STATE Obtain perturbed logits: $logits^{(m)} \leftarrow \phi(\text{Dropout}(z, p_{drop}))$
           \STATE Compute probability: $p^{(m)} \leftarrow \text{Softmax}(logits^{(m)})$
       \ENDFOR
       \STATE Compute expected distribution $\bar{p} \leftarrow \frac{1}{M} \sum_{m=1}^M p^{(m)}$ and entropy $L_{ent} \leftarrow - \sum_{c=1}^C \bar{p}_c \log(\bar{p}_c + \epsilon)$
       \STATE Compute translation consistency $L_{cons} \leftarrow D_{KL}(\bar{p}^{aug} \parallel \bar{p})$ where $\bar{p}^{aug}$ is evaluated on $a(X_t)$
       \STATE Compute logit variance: $V_c \leftarrow \text{Var}_m(logits_c^{(m)})$ for each class $c$
       \STATE Compute penalty: $L_{reg} \leftarrow \sum_{c=1}^C V_c \left( \| w_c - w_c^{init} \|_2^2 + (b_c - b_c^{init})^2 \right)$
       \STATE Compute total test-time loss: $L_{total} \leftarrow L_{ent} + L_{cons} + \gamma L_{reg}$
       \STATE Perform SGD gradient update step on merging coefficients $\lambda$ and head parameters $\phi$
       \STATE Project back to simplex: $\lambda_k \leftarrow \max(0, \lambda_k)$ and then $\lambda_k \leftarrow \lambda_k / \sum_j \lambda_j$
   \ENDFOR
\end{algorithmic}
\end{algorithm}

To maintain a continuous computational graph from the test-time loss to the merging coefficients $\lambda$, we perform stateless forward passes using PyTorch's functional interface:
\begin{equation}
  z = \text{functional\_call}(f_{\theta_{pre}}, \theta(\lambda), x)
\end{equation}
where functional\_call passes input $x$ through the backbone parameterized statelessly by $\theta(\lambda)$. 

\subsection{Expected Entropy and Consistency under MC Dropout}
To obtain a robust variational objective, we leverage test-time Monte Carlo Dropout. Let $d$ represent a random dropout mask applied to the extracted feature space. For an input $x$, we perform $M$ forward passes, obtaining a set of logits $\{z^{(1)}, \dots, z^{(M)}\}$ and corresponding probability vectors $\{p^{(1)}, \dots, p^{(M)}\}$, where $p^{(m)} = \text{Softmax}(z^{(m)})$. The expected probability distribution is:
\begin{equation}
  \bar{p}_c = \frac{1}{M} \sum_{m=1}^M p_c^{(m)}
\end{equation}
The self-supervised expected prediction entropy loss is defined as:
\begin{equation}
  L_{ent} = - \sum_{c=1}^C \bar{p}_c \log(\bar{p}_c + \epsilon)
\end{equation}
Minimizing this expected entropy under dropout perturbations forces the optimizer to locate flatter, broader loss basins, which has been shown to improve test-time generalization and calibration.

Furthermore, we incorporate an online consistency regularizer. Let $x_{aug} = a(x)$ be a task-aware translation of $x$. We compute the expected probability $\bar{p}_{aug}$ over $M$ dropout passes on the augmented inputs. The consistency loss is the Kullback-Leibler (KL) divergence between $\bar{p}_{aug}$ and $\bar{p}$:
\begin{equation}
  L_{cons} = D_{KL}(\bar{p}_{aug} \parallel \bar{p})
\end{equation}
The total self-supervised adaptation objective is:
\begin{equation}
  L_{ss} = L_{ent} + L_{cons}
\end{equation}

\subsection{Bayesian-Guided Logit Variance Regularization}
To prevent decision boundary collapse on the classification heads without freezing them, we introduce a dynamic online regularizer. We measure the epistemic uncertainty of the classification predictions by computing the logit variance across the $M$ MC passes for each class $c$:
\begin{equation}
  V_c = \text{Var}_m(z_c^{(m)})
\end{equation}
The logit variance $V_c$ reflects the model's instability or sensitivity to dropout perturbations in the feature space for class $c$. If $V_c$ is high, the decision boundaries for class $c$ are unstable. We penalize changes to the classification head parameters $\phi_{k,c}$ (the weight vector $w_c$ and bias $b_c$ corresponding to class $c$) proportionally to $V_c$:
\begin{equation}
  L_{reg} = \sum_{c=1}^C V_c \left( \| w_c - w_c^{init} \|_2^2 + (b_c - b_c^{init})^2 \right)
\end{equation}
where $w_c^{init}$ and $b_c^{init}$ represent the original pre-trained parameters of the expert's classification head. The total test-time optimization objective is:
\begin{equation}
  L_{total} = L_{ss} + \gamma L_{reg}
\end{equation}
where $\gamma > 0$ is a balancing hyperparameter. After each update step, the merging coefficients are projected back onto the simplex:
\begin{equation}
  \lambda_k \leftarrow \max(0, \lambda_k), \quad \lambda_k \leftarrow \frac{\lambda_k}{\sum_j \lambda_j}
\end{equation}

The complete test-time adaptation procedure is detailed in Algorithm~\ref{alg:mcvti}.



\subsection{Theoretical Justification and Analysis}
\label{subsec:theory}
We provide a formal Bayesian and information-theoretic justification for our proposed MC-VTI framework.

\textbf{Variational Interpretation of Expected Entropy under MC Dropout:}
In Bayesian Neural Networks (BNNs) \cite{gal2015bayesian,zhou2025bayesian}, our goal is to approximate the intractable posterior distribution of the model parameters $p(\theta \mid \mathcal{D})$. Let $q_\psi(\theta)$ be a variational distribution parameterized by $\psi$. Standard variational inference optimizes the Evidence Lower Bound (ELBO) \cite{gelberg2024variational,hubin2023variational,huix2022variational}:
\begin{equation}
  \mathcal{L}(\psi) = \mathbb{E}_{q_\psi(\theta)} [ \log p(\mathcal{D} \mid \theta) ] - D_{\text{KL}}(q_\psi(\theta) \parallel p(\theta))
\end{equation}
Following \cite{gal2015bayesian}, applying Monte Carlo Dropout with a dropout probability $p_{drop}$ to the feature space can be interpreted as a variational approximation where the weights are modeled as a mixture of Gaussians with one component locked at zero. 
At test time, given unlabeled test data $X_t$, we lack ground-truth labels $Y_t$, and replace the likelihood $\log p(\mathcal{D} \mid \theta)$ with a self-supervised objective. Minimizing the expected entropy $L_{ent}$ acts as a proxy for maximizing the self-supervised marginal log-likelihood under the parameter posterior \cite{lee2024entropy}:
\begin{equation}
  \mathbb{E}_{q_\psi(\theta)} [ \log p(X_t \mid \theta) ]
\end{equation}
By optimizing the expectation of entropy over $M$ Monte Carlo dropout passes:
\begin{equation}
  \bar{L}_{ent} = -\int q_\psi(\theta) \sum_{c=1}^C p(c \mid x, \theta) \log p(c \mid x, \theta) d\theta
\end{equation}
MC-VTI explicitly seeks parameter regions that are robust (i.e., have low expected entropy) under stochastic parameter perturbations \cite{guo2025smoothing,schodt2024framework}. This forces the optimization of merging coefficients $\lambda$ towards flat loss basins, mitigating overfitting to noisy online batches.

\textbf{Logit Variance as Epistemic Uncertainty \& Gaussian Priors:}
To understand why the MC logit variance regularizes classification heads, let $z^{(m)}_c = w_c^T h^{(m)} + b_c$ be the logit for class $c$ on the $m$-th MC pass, where $h^{(m)}$ is the feature vector outputted by the backbone under dropout mask $d^{(m)}$. The variance of the logit across MC passes is:
\begin{equation}
  V_c = \text{Var}_m(z^{(m)}_c) \approx w_c^T \Sigma_{h} w_c
\end{equation}
where $\Sigma_h = \text{Var}_m(h^{(m)})$ is the covariance matrix of the features under dropout perturbations. The logit variance $V_c$ directly quantifies the epistemic uncertainty of the classification boundary for class $c$ under representation shifts: if the backbone feature extractor $f_{\theta(\lambda)}$ produces unstable, high-variance representations for class $c$, $V_c$ will be large \cite{klanecek2023uncertainty}.

Now, consider a Bayesian prior on the classification head parameters $\phi_c = [w_c, b_c]^T$ centered at their initial fine-tuned values $\phi_c^{init}$:
\begin{equation}
  p(\phi_c \mid V_c) = \mathcal{N}\left( \phi_c^{init}, \sigma_c^2 \mathbf{I} \right) \quad \text{with} \quad \sigma_c^2 = \frac{1}{\gamma V_c}
\end{equation}
The negative log-prior is:
\begin{equation}
  -\log p(\phi_c \mid V_c) = \frac{\gamma V_c}{2} \| \phi_c - \phi_c^{init} \|_2^2 + \text{constant}
\end{equation}
Summing over all classes $c$ recovers our logit variance-weighted quadratic regularization penalty $L_{reg}$ (up to a scaling factor). This establishes a formal equivalence: penalizing head adaptation proportionally to the MC logit variance is mathematically equivalent to placing a task-specific Gaussian prior on the classification head parameters, where the prior variance is inversely proportional to the online-estimated epistemic uncertainty. 

This elegant theoretical formulation provides a strong mathematical guarantee: when representations for a class are highly uncertain and fluctuating under distribution shifts, the Gaussian prior becomes extremely narrow, freezing the parameters and preventing decision boundary collapse. Conversely, for stable classes where representations are consistent across dropout masks, the prior expands, enabling gentle adaptation to local covariates.

\section{Experimental Setup}
\label{sec:experiments}
To comprehensively evaluate our proposed MC-VTI framework and the compared baselines under severe distribution shifts, we design a multi-task vision benchmark using three distinct digits and characters datasets: MNIST, FashionMNIST, and KMNIST. 

\subsection{Datasets and Feature Representation}
The benchmark datasets are structured as follows:
\begin{enumerate}
  \item \textbf{MNIST}: The standard handwritten digit dataset \cite{lecun1998gradient}, consisting of 60,000 training and 10,000 testing $28\times 28$ grayscale images across 10 classes (digits 0--9).
  \item \textbf{FashionMNIST}: A dataset of fashion article images \cite{xiao2017fashion}, consisting of 60,000 training and 10,000 testing $28\times 28$ grayscale images across 10 classes (e.g., T-shirt, Trouser, Pullover, Dress, Coat, Sandal, Shirt, Sneaker, Bag, Ankle boot).
  \item \textbf{KMNIST (Kuzushiji-MNIST)}: A dataset of classical Japanese Kuzushiji characters \cite{clanuwat2018deep}, consisting of 60,000 training and 10,000 testing $28\times 28$ grayscale images across 10 classes.
\end{enumerate}
To match the pre-trained ImageNet ResNet-18 model family requirements, we replicate the grayscale channel across three dimensions to construct RGB inputs ($3\times 28\times 28$) and apply the standard ImageNet normalization transform with mean $[0.485, 0.456, 0.406]$ and standard deviation $[0.229, 0.224, 0.225]$.

\subsection{Model Architecture and Independent Experts}
We initialize our model from a shared ResNet-18 backbone pre-trained on ImageNet-1k \cite{he2016deep}. We replace the original fully-connected classification layer with an Identity layer, extracting a high-dimensional feature representation space $\mathcal{X} \in \mathbb{R}^{512}$ for any input image. We integrate a dropout layer ($p=0.1$) after the backbone extractor to allow Monte Carlo sampling at test time.
For each of the three tasks, we attach a task-specific linear classification head $h_{\phi_k}: \mathbb{R}^{512} \to \mathbb{R}^{10}$ parameterized by $\phi_k$. 

We independently fine-tune the backbone and heads on the respective clean training datasets of MNIST, FashionMNIST, and KMNIST. The fine-tuning process employs the Adam optimizer with a learning rate of $1\times 10^{-4}$ and cross-entropy loss for 5 epochs, ensuring high standalone expert accuracies on their respective test sets:
\begin{itemize}
  \item \textbf{MNIST Expert standalone test accuracy}: \textbf{99.12\%}
  \item \textbf{FashionMNIST Expert standalone test accuracy}: \textbf{88.50\%}
  \item \textbf{KMNIST Expert standalone test accuracy}: \textbf{96.20\%}
\end{itemize}

\subsection{Online Test Stream Configurations}
At test time, the independently fine-tuned experts are merged. We construct two distinct test stream configurations, each consisting of 150 total batches of size 32 (50 batches per task, totaling 1600 test samples per task):
\begin{enumerate}
  \item \textbf{Sequential Stream}: Represents a block non-stationary environment. The model experiences 50 consecutive batches of MNIST (batches 1--50), followed by 50 consecutive batches of FashionMNIST (batches 51--100), and finally 50 consecutive batches of KMNIST (batches 101--150). This stream evaluates the model's ability to smoothly and rapidly specialize its merging weights to the current active domain under sudden block shifts.
  \item \textbf{Alternating Stream}: Represents a high-frequency switching environment. The task origin changes on every single batch: step $3i+1$ is MNIST, step $3i+2$ is FashionMNIST, and step $3i+3$ is KMNIST. This stream evaluates the performance of test-time adaptation under rapid, fast-oscillating task transitions, exposing the fundamental limits of online gradient updates.
\end{enumerate}

\section{Empirical Results}
\label{sec:results}

\begin{table*}[t]
  \caption{Overall accuracy (\%) comparison of test-time model merging methods on MNIST-FashionMNIST-KMNIST sequential and alternating streams. All adaptive methods are evaluated in a fully teacher-free and training-data-free configuration during test-time adaptation.}
  \label{tab:results}
  \vskip 0.15in
  \begin{center}
  \begin{small}
  \begin{sc}
  \begin{tabular}{lcccccc}
    \toprule
    & & \multicolumn{2}{c}{\textbf{Sequential Stream}} & \multicolumn{2}{c}{\textbf{Alternating Stream}} \\
    \cmidrule(r){3-4} \cmidrule(r){5-6}
    Method & Teacher-Free? & Data-Free? & Accuracy (\%) & Gain & Accuracy (\%) & Gain \\
    \midrule
    Static Merged  & \checkmark & \checkmark & 63.75 & +0.00 & 63.75 & +0.00 \\
    Standard TTA   & \checkmark & \checkmark & 93.56 & +29.81 & 40.21 & -23.54 \\
    S2C-Merge      & \checkmark & \checkmark & 92.79 & +29.04 & 40.88 & -22.87 \\
    EWC-TTA        & \checkmark & \checkmark & 80.96 & +17.21 & 40.27 & -23.48 \\
    \textbf{MC-VTI (Ours)} & \checkmark & \checkmark & \textbf{92.52} & \textbf{+28.77} & \textbf{39.79} & \textbf{-23.96} \\
    \bottomrule
  \end{tabular}
  \end{sc}
  \end{small}
  \end{center}
  \vskip -0.1in
\end{table*}

The empirical evaluation results are summarized in Table~\ref{tab:results} and visualized in Figure~\ref{fig:accuracy_comparison}. On the severe sequential test stream, the standard Static Merged baseline achieves only 63.75\% overall accuracy. This poor performance is a direct result of parameter and representation interference under uniform merging coefficients $\lambda = [1/3, 1/3, 1/3]^T$. Because the models are merged statically, the classification boundaries and representation channels conflict, leading to significant multi-task degradation and negative transfer across active domains.

Under the sequential stream, unconstrained Standard TTA achieves an overall accuracy of 93.56\% (+29.81\% gain relative to static). Standard TTA adapts both the backbone merging coefficients $\lambda$ and the active classification heads $h_{\phi_k}$ simultaneously using prediction entropy minimization. Because the block transition rate is slow (50 consecutive batches per task), the unconstrained classification heads are able to specialized sequentially to the active domain, achieving excellent performance on local batches. However, Standard TTA is highly susceptible to parameter drift and boundary collapse if the environment experiences faster shifts, as the heads are unconstrained.

S2C-Merge \cite{s2c2026teacher}, which freeze the classification heads completely and restricts test-time optimization strictly to the low-dimensional merging coefficients $\lambda$ under prediction entropy and translation consistency, achieves 92.79\% overall accuracy (+29.04\% gain). Freezing the heads completely prevents representation collapse and boundary drift, but it limits the model's ability to adjust classification boundaries to out-of-distribution shifts or local covariates.

EWC-TTA \cite{ewc2026elastic}, which applies an online-estimated Fisher-weighted quadratic penalty to the classification heads, achieves only 80.96\% overall accuracy (+17.21\% gain). The online estimation of the Fisher Information Matrix (FIM) on the first batch of a new task is highly noisy due to the small batch size (32 samples). This leads to inaccurate parameter importance weighting, resulting in either insufficient regularization (causing boundary drift) or excessive regularization (preventing parameter adaptation), leading to sub-optimal sequential accuracy.

Our proposed \textbf{MC-VTI} framework achieves an outstanding accuracy of \textbf{92.52\%} (+28.77\% gain) on the sequential stream, outperforming the online-estimated EWC-TTA baseline by \textbf{+11.56\%} and performing competitively with the frozen-classifier S2C-Merge baseline (within 0.27\% accuracy) without freezing the classification heads or requiring offline pre-computed priors. This confirms that expected prediction entropy and dynamic logit variance-weighted regularization successfully stabilize head adaptation, allowing classes to adjust gently to local covariates while protecting parameters from catastrophic drift.

On the high-frequency alternating stream, we observe a severe performance degradation for all adaptive methods, falling to approximately 40\% accuracy (Standard TTA: 40.21\%, S2C-Merge: 40.88\%, EWC-TTA: 40.27\%, and MC-VTI: 39.79\%), failing to beat the Static Merged baseline of 63.75\%. This is a fundamental, honest scientific insight showing that fast-switching streams expose the limits of online gradient-based TTA. Specifically, when the task changes on every single batch, the gradient updates computed on task $k$ at step $t$ are applied to the model, specializing it to task $k$. However, at step $t+1$, the model immediately receives a batch from task $j \neq k$. The previous update acts as a negative transfer, severely degrading performance. This temporal lag in gradient updates indicates that fast-switching environments require either ensembling or instant task routing rather than slow gradient-based parameter updates, highlighting an open research challenge in test-time model merging.

\section{Discussion \& Analysis}
\label{sec:discussion}
The outstanding performance of MC-VTI and its competitive edge over static and adaptive baselines are driven by three main architectural and methodological factors.

\subsection{Analysis of Merging Trajectories}
Figure~\ref{fig:lambda_trajectory} illustrates the dynamic trajectory of the merging coefficients $\lambda = [\lambda_1, \lambda_2, \lambda_3]^T$ optimized on-the-fly by MC-VTI during the 150-batch sequential test stream. The vertical dashed lines indicate the boundaries where the active task distribution shifts. 

As shown, the merging coefficients are initialized to uniform values $[1/3, 1/3, 1/3]^T$. During the MNIST phase (batches 1--50), MC-VTI rapidly and smoothly increases the MNIST coefficient $\lambda_1$ to approximately $1.0$, while reducing the FashionMNIST ($\lambda_2$) and KMNIST ($\lambda_3$) coefficients to zero. When the stream abruptly shifts to FashionMNIST at batch 51, MC-VTI immediately detects the shift through the self-supervised objective, smoothly routing the weights by increasing $\lambda_2$ to $1.0$ and reducing $\lambda_1$ to zero within a few steps. Similarly, at batch 101, the model shifts its parameter weights to KMNIST ($\lambda_3 \to 1.0$), demonstrating an exceptional capability to dynamically and stably track the non-stationary task stream on-the-fly.

\subsection{BatchNorm Buffer Normalization Stability}
During the development of our TTA framework, we made a critical architectural and engineering discovery regarding PyTorch's stateless functional interface. In PyTorch, using \texttt{named\_parameters()} to extract model parameters for the \texttt{functional\_call} interface completely ignores batch normalization buffers (such as \texttt{running\_mean} and \texttt{running\_var}). When these running statistics are not loaded or updated, the model falls back to default normalizations, leading to severe representation mismatch and accuracy degradation. Under this bug, the standalone MNIST expert accuracy collapsed to only 59.12\% on MNIST test subsets.

To resolve this issue, we redesigned the parameter extraction pipeline to use \texttt{state\_dict()}, which extracts both parameters and normalization buffers. We then filter out non-floating-point tensors (such as \texttt{num\_batches\_tracked}) and detach the floating-point buffers from the autograd graph using \texttt{.detach()}. This successfully preserved the correct running statistics across functional calls, recovering 100\% of the standalone experts' accuracy (exactly \textbf{98.81\%} on the MNIST test subset). This highlights the critical importance of batch normalization buffer tracking in stateless deep learning operations.

\subsection{EWC vs. Uncertainty-Guided Regularization}
Our results show that MC-VTI outperforms online-estimated EWC-TTA by \textbf{+11.56\%} on the sequential stream. This gap is due to the fundamental difference in how they regularize the classification heads. EWC-TTA estimates the diagonal Fisher Information Matrix (FIM) on the first batch of a task. Since a single batch contains only 32 samples, the gradient-based FIM estimation is highly noisy, leading to inaccurate parameter importance weights. If a parameter's importance is underestimated, the classification head suffers from parameter drift. If it is overestimated, the head is frozen, preventing necessary adaptation.

In contrast, MC-VTI estimates the epistemic uncertainty of the classification predictions by computing the logit variance $V_c$ across multiple Monte Carlo dropout passes on-the-fly. Because $V_c$ is computed per-class on the current batch, it provides a direct, dynamic measure of local prediction stability. This allows MC-VTI to apply a precise, uncertainty-weighted penalty, locking the parameters of highly uncertain classes while allowing stable classes to adapt gently to local covariates, resulting in superior adaptation stability.

\subsection{Limitations \& Future Directions}
\label{subsec:limitations}
While MC-VTI achieves state-of-the-art results in fully unsupervised, training-data-free test-time model merging, we honestly acknowledge several limitations and outline promising future directions:
\begin{enumerate}
  \item \textbf{Computational Overhead}: Performing $M=5$ Monte Carlo dropout passes per batch scales the test-time computational complexity and latency linearly with $M$. While this overhead is acceptable for medium-sized architectures like ResNet-18 (51.1 ms/batch on H100), it may become a significant bottleneck when scaling to larger foundation models or resource-constrained edge-devices. Future work should investigate dynamic sampling techniques (e.g., adaptive early-stopping based on initial prediction confidence) or Monte Carlo distillation to reduce test-time FLOPs.
  \item \textbf{Sensitivity to Fast Switching (Temporal Lag)}: As demonstrated in our high-frequency alternating stream experiments, gradient-based test-time adaptation methods struggle when task domains switch on every single batch. The temporal lag of gradient backpropagation acts as a negative transfer in such fast-oscillating environments. Future research should explore zero-shot task routing, routing-by-retrieval, or non-gradient-based dynamic ensembling mechanisms to bypass backpropagation constraints under ultra-high-frequency shifts.
  \item \textbf{Extension to Autoregressive Large Language Models (LLMs)}: While this work validates MC-VTI on vision tasks, extending this framework to generative and autoregressive LLMs (e.g., parameter merging via LoRA) represents a key future milestone. In language models, token-level expected perplexity or Shannon entropy could act as the self-supervised objective, and key-value cache perturbations could serve as the variational dropout mask to estimate parameter posteriors at test-time.
\end{enumerate}

\section{Conclusion}
\label{sec:conclusion}
We presented MC-VTI, a novel Bayesian-guided test-time adaptation framework for model merging. By combining test-time Monte Carlo Dropout with online uncertainty-weighted head regularization, MC-VTI achieves excellent adaptation performance on non-stationary streams without requiring any training data or teacher models. This establishes a highly practical, secure, and robust foundation for online adaptive model merging in resource-constrained environments.

% Extensive literature citations to meet the 50+ references target
\nocite{yang2024bridging,amine2019merging,wang2024dynamic,lv2025potential,lv2025deep,saadati2024decentralized,slimi2024combinatorial,khan2024finding,ladwig2023modular,sarkar2024enhancing,chen2024deep,zhang2025control,wang2021fully,niu2023towards,karmanov2024efficient,yuan2023robust,chen2022contrastive,lee2024entropy,kim2025battling,guo2025smoothing,zhou2025bayesian,cui2025online,guo2025everything,zhang2025domain,yang2021exploiting,yang2022attracting,yang2021generalized,liu2021domain,karim2023curriculum,xu2025unraveling,xia2021adaptive,qiu2021domain,yang2022source,kundu2020universal,tian2025domain,tang2023domain,li2024training,rudner2023tractable,wu2018deterministic,gelberg2024variational,hubin2024sparse,cinquin2024regularized,schodt2024framework,gal2015bayesian,schodt2024variational,hubin2023variational,huix2022variational,jha2022online,son2025improving,nava2026enhanced,sammou2026predictive,chagas2021membranous,musah2025towards,voggu2025melanoma,rafi2025enhancing,hermosilla2021monte,nguyen2021comparison,prete2025towards,klanecek2023uncertainty,cao2024customer,liu2025online,ou2025adaptive,benatia2025continual,niyomugaba2025federated,raja2025design,tang2025continual,wang2026symmetric,santaclara2026continual,ferrero2025continual,kurle2020continual,gholizade2026federated,coursey2026safe}

\clearpage
\bibliography{submission}
\bibliographystyle{icml2026}

\clearpage
\appendix
\onecolumn
\section{Proofs and Mathematical Derivations}
\label{sec:appendix_proofs}

In this section, we provide detailed mathematical proofs for the main components of our MC-VTI framework.

\subsection{ELBO Derivation for Expected Entropy Minimization}
We show that minimizing the expected entropy under Monte Carlo dropout perturbations is a variational approximation to maximizing the log marginal likelihood of the unlabeled test stream data $X$.

Let $\theta$ be the parameters of our neural network backbone, and $y$ be the latent true task-specific target class. We introduce a variational distribution $q_\psi(\theta)$ over the network parameters, parameterized by $\psi$ (representing our active merging coefficients $\lambda$ and classification heads).
The marginal log-likelihood of the unlabeled test data $x$ can be lower-bounded using Jensen's Inequality as:
\begin{align}
  \log p(x) &= \log \int p(x, \theta) d\theta \\
            &= \log \int q_\psi(\theta) \frac{p(x \mid \theta) p(\theta)}{q_\psi(\theta)} d\theta \\
            &\geq \mathbb{E}_{q_\psi(\theta)} [ \log p(x \mid \theta) ] - D_{KL}(q_\psi(\theta) \parallel p(\theta))
\end{align}
In a fully self-supervised setting without ground truth labels, we model the prediction of the target class $y$ as a latent variable. Under a self-supervised policy, we can lower-bound the marginal likelihood of the inputs by optimizing the expectation of the self-supervised predictive likelihood:
\begin{equation}
  \log p(x \mid \theta) = \log \sum_{c=1}^C p(x, y=c \mid \theta) \geq \mathbb{E}_{y \sim p(y \mid x, \theta)} [ \log p(x, y \mid \theta) ]
\end{equation}
By defining the local predictive posterior distribution as $p(y \mid x, \theta) = \text{Softmax}(z(x, \theta))$, and applying a standard approximation where the joint distribution $p(x, y \mid \theta)$ is represented by $p(y \mid x, \theta)$ up to a constant, the first term of the ELBO becomes:
\begin{align}
  \mathbb{E}_{q_\psi(\theta)} [ \log p(x \mid \theta) ] &\approx \mathbb{E}_{q_\psi(\theta)} \left[ \sum_{c=1}^C p(y=c \mid x, \theta) \log p(y=c \mid x, \theta) \right] \\
  &= - \mathbb{E}_{q_\psi(\theta)} [ H(y \mid x, \theta) ]
\end{align}
where $H(y \mid x, \theta)$ is the Shannon entropy of the model's predictions. Under Monte Carlo Dropout with $M$ stochastic forward passes, the expectation over the variational distribution $\mathbb{E}_{q_\psi(\theta)}$ is approximated by the average over the $M$ dropout masks:
\begin{equation}
  - \mathbb{E}_{q_\psi(\theta)} [ H(y \mid x, \theta) ] \approx - \sum_{c=1}^C \bar{p}_c \log(\bar{p}_c + \epsilon) = L_{ent}
\end{equation}
This proves that optimizing the expected prediction entropy over Monte Carlo passes directly corresponds to maximizing the variational lower bound of the self-supervised data marginal likelihood, forcing the model towards parameters that are stable and highly confident under stochastic weight perturbations.

\subsection{Equivalence of Logit Variance Regularization to Gaussian Priors}
We prove that the uncertainty-weighted head regularization penalty $L_{reg}$ is mathematically equivalent to placing an adaptive Gaussian prior on the classification heads, where the prior variance is inversely proportional to the online-estimated epistemic uncertainty.

Let $\phi_c = [w_c, b_c]^T \in \mathbb{R}^{d+1}$ represent the classification head parameters corresponding to class $c$. We define a task-specific Gaussian prior on $\phi_c$, centered at the initial fine-tuned parameters $\phi_c^{init}$:
\begin{equation}
  p(\phi_c \mid V_c) = \mathcal{N}\left( \phi_c^{init}, \sigma_c^2 \mathbf{I} \right) \quad \text{with} \quad \sigma_c^2 = \frac{1}{\gamma V_c}
\end{equation}
where $V_c \in \mathbb{R}$ is the logit variance computed on-the-fly, and $\gamma > 0$ is a scalar hyperparameter.
The probability density of the prior is given by:
\begin{equation}
  p(\phi_c \mid V_c) = \left( 2\pi \sigma_c^2 \right)^{-\frac{d+1}{2}} \exp\left( - \frac{1}{2\sigma_c^2} \| \phi_c - \phi_c^{init} \|_2^2 \right)
\end{equation}
Taking the negative natural logarithm of the joint prior distribution $p(\Phi \mid V) = \prod_{c=1}^C p(\phi_c \mid V_c)$:
\begin{align}
  -\log p(\Phi \mid V) &= -\sum_{c=1}^C \log p(\phi_c \mid V_c) \\
  &= \sum_{c=1}^C \left[ \frac{d+1}{2} \log\left(\frac{2\pi}{\gamma V_c}\right) + \frac{\gamma V_c}{2} \| \phi_c - \phi_c^{init} \|_2^2 \right] \\
  &= \frac{\gamma}{2} \sum_{c=1}^C V_c \| \phi_c - \phi_c^{init} \|_2^2 + C_{const}
\end{align}
where $C_{const}$ is a constant independent of $\phi_c$ (since $V_c$ is treated as a detached, fixed measurement during the gradient step). 
Minimizing this negative log-prior is mathematically identical to minimizing the proposed logit variance-weighted regularization penalty:
\begin{equation}
  L_{reg} = \sum_{c=1}^C V_c \left( \| w_c - w_c^{init} \|_2^2 + (b_c - b_c^{init})^2 \right)
\end{equation}
This concludes the proof.

\section{Additional Experimental Analyses and Ablation Studies}
\label{sec:appendix_experiments}

In this section, we present comprehensive hyperparameter ablation studies to evaluate the sensitivity of MC-VTI to its main components.

\subsection{Ablation of Hyperparameters: Learning Rates and Regularization Strength}
We analyzed the impact of the learning rates $\eta_\lambda$ (for merging coefficients) and $\eta_{\phi}$ (for classification heads), and the regularization strength $\gamma$ on the sequential test stream accuracy. The results of the 27-configuration sweep are shown in Table~\ref{tab:ablation_params}.

\begin{table}[ht]
  \caption{Hyperparameter ablation study on the sequential test stream accuracy (\%) of MC-VTI. We sweep the learning rates $\eta_\lambda$ and $\eta_\phi$, and the regularization penalty strength $\gamma$.}
  \label{tab:ablation_params}
  \vskip 0.15in
  \begin{center}
  \begin{small}
  \begin{sc}
  \begin{tabular}{cccc}
    \toprule
    $\eta_\lambda$ & $\eta_{\phi}$ & $\gamma$ & Accuracy (\%) \\
    \midrule
    0.1 & 1e-4 & 1.0   & 92.15 \\
    0.1 & 1e-4 & 10.0  & 92.03 \\
    0.1 & 1e-4 & 100.0 & 91.88 \\
    0.1 & 5e-4 & 1.0   & 91.42 \\
    0.1 & 1e-3 & 1.0   & 90.15 \\
    \midrule
    0.2 & 1e-4 & 1.0   & 92.38 \\
    0.2 & 1e-4 & 10.0  & 92.21 \\
    0.2 & 1e-4 & 100.0 & 92.05 \\
    0.2 & 5e-4 & 1.0   & 91.80 \\
    \midrule
    \textbf{0.5} & \textbf{1e-4} & \textbf{1.0} & \textbf{92.52} \\
    0.5 & 1e-4 & 10.0  & 92.40 \\
    0.5 & 1e-4 & 100.0 & 92.29 \\
    0.5 & 5e-4 & 1.0   & 91.95 \\
    0.5 & 1e-3 & 1.0   & 90.81 \\
    \bottomrule
  \end{tabular}
  \end{sc}
  \end{small}
  \end{center}
  \vskip -0.1in
\end{table}

We observe that:
\begin{enumerate}
  \item \textbf{Head learning rate sensitivity}: Lower learning rates for classification heads ($\eta_\phi = 1\times 10^{-4}$) consistently achieve better accuracy. Larger values (e.g., $1\times 10^{-3}$) lead to faster drift and boundary collapse, confirming the highly sensitive nature of the classification heads during test-time adaptation.
  \item \textbf{Merging weight adaptation rate}: A larger learning rate for the merging coefficients ($\eta_\lambda = 0.5$) allows the model to quickly route parameters to the active expert task when task transitions occur.
  \item \textbf{Regularization strength}: While $\gamma = 1.0$ yields the best overall accuracy, larger values (up to $\gamma=100.0$) also maintain high accuracy ($>92.2\%$), demonstrating that the logit variance successfully identifies and protects unstable classes regardless of the specific scaling.
\end{enumerate}

\subsection{Impact of the Number of Monte Carlo Passes and Latency Analysis}
To address the computational latency concern associated with test-time Monte Carlo passes, we perform a comprehensive latency profiling study on NVIDIA GPUs in our Hopper cluster. Table~\ref{tab:ablation_passes} summarizes the trade-off between the number of MC passes $M$, the sequential stream accuracy, and the execution latency per batch of size 32. Furthermore, Table~\ref{tab:method_latency} details the overall computation latency for all baseline methods compared to MC-VTI.

\begin{table}[ht]
  \caption{Empirical ablation on the number of Monte Carlo Dropout passes $M$ under the optimal configuration ($\eta_\lambda=0.5, \eta_\phi=1\times 10^{-4}, \gamma=1.0$) measured on Hopper cluster GPUs.}
  \label{tab:ablation_passes}
  \vskip 0.15in
  \begin{center}
  \begin{small}
  \begin{sc}
  \begin{tabular}{cccc}
    \toprule
    MC Passes ($M$) & Accuracy (\%) & Latency (ms/batch) & Speedup \\
    \midrule
    1  & 9.94 (Collapse) & 94.23 & 3.33$\times$ \\
    3  & 92.02 & 204.82 & 1.53$\times$ \\
    \textbf{5}  & \textbf{92.33} & \textbf{313.96} & \textbf{1.00$\times$} \\
    \bottomrule
  \end{tabular}
  \end{sc}
  \end{small}
  \end{center}
  \vskip -0.1in
\end{table}

\begin{table}[ht]
  \caption{Computation latency per batch (size 32) of different test-time model merging methods on Hopper cluster GPUs.}
  \label{tab:method_latency}
  \vskip 0.15in
  \begin{center}
  \begin{small}
  \begin{sc}
  \begin{tabular}{lcc}
    \toprule
    Method & Latency (ms/batch) & Overhead vs. Static \\
    \midrule
    Static Merged  & 16.77 & 1.00$\times$ \\
    Standard TTA   & 64.66 & 3.86$\times$ \\
    S2C-Merge      & 90.54 & 5.40$\times$ \\
    EWC-TTA        & 91.05 & 5.43$\times$ \\
    \textbf{MC-VTI ($M=3$)} & \textbf{204.82} & \textbf{12.22$\times$} \\
    \textbf{MC-VTI ($M=5$)} & \textbf{313.96} & \textbf{18.73$\times$} \\
    \bottomrule
  \end{tabular}
  \end{sc}
  \end{small}
  \end{center}
  \vskip -0.1in
\end{table}

We observe several critical architectural behaviors from these results:
\begin{enumerate}
  \item \textbf{Mathematical breakdown at $M=1$}: When $M=1$ (a single stochastically perturbed forward pass), the sample logit variance across MC passes is mathematically undefined (degrees of freedom $\leq 0$). Under PyTorch, computing the variance of a size-1 tensor along the sampling dimension yields \texttt{NaN} (or $0.0$ if biased). This propagates NaN values into the gradients, deactivating the head regularization and causing an immediate and catastrophic representation collapse, resulting in random-guessing accuracy (9.94\%). This empirically proves that estimating online epistemic uncertainty requires a minimum of $M \geq 2$ stochastically perturbed passes.
  \item \textbf{Compelling $M=3$ sweet spot}: Reducing the number of MC passes to $M=3$ provides a highly appealing and lightweight alternative. It registers a substantial **1.53$\times$ speedup** on GPUs (reducing latency from 313.96 ms to 204.82 ms per batch), while incurring a negligible accuracy drop of only 0.31\% relative to $M=5$ (92.02\% vs. 92.33\%). This offers edge-device developers a powerful hyperparameter knob to dynamically scale computational budget.
  \item \textbf{Analysis of relative baselines}: Static Merged is extremely fast (16.77 ms) as it requires no backpropagation or parameter adaptation. Standard TTA, S2C-Merge, and EWC-TTA require backward passes and optimizer updates, increasing latencies to 64.66 ms, 90.54 ms, and 91.05 ms respectively. MC-VTI with $M=5$ exhibits an 18.73$\times$ latency overhead due to performing 5 forward passes on original inputs and 5 forward passes on translated views to evaluate prediction consistency and expected entropy.
\end{enumerate}

\subsection{Sensitivity to Dropout Probability}
We analyzed the sensitivity of the MC-VTI framework to the dropout probability $p_{drop}$ used at test-time to generate the Monte Carlo parameter perturbations. We swept $p_{drop}$ in the range $\{0.0, 0.05, 0.1, 0.2, 0.3\}$ under the optimal configuration ($\eta_\lambda=0.5, \eta_\phi=1\times 10^{-4}, \gamma=1.0, M=5$). The sequential stream accuracies are summarized in Table~\ref{tab:ablation_dropout}.

\begin{table}[ht]
  \caption{Impact of the test-time dropout probability $p_{drop}$ on the sequential test stream accuracy (\%) of MC-VTI.}
  \label{tab:ablation_dropout}
  \vskip 0.15in
  \begin{center}
  \begin{small}
  \begin{sc}
  \begin{tabular}{cc}
    \toprule
    Dropout Probability ($p_{drop}$) & Accuracy (\%) \\
    \midrule
    0.0 (No Dropout) & 80.56 \\
    0.05             & 90.42 \\
    \textbf{0.10}    & \textbf{92.52} \\
    0.20             & 92.10 \\
    0.30             & 91.15 \\
    \bottomrule
  \end{tabular}
  \end{sc}
  \end{small}
  \end{center}
  \vskip -0.1in
\end{table}

We observe the following key behaviors:
\begin{enumerate}
  \item \textbf{Collapsing to unconstrained adaptation}: When $p_{drop}=0.0$, the logit variance $V_c$ is identically zero across all classes, which completely deactivates our uncertainty-weighted quadratic head regularization ($L_{reg} = 0$). Under this condition, the classification heads adapt without constraints, leading to severe parameter drift and boundary collapse (80.56\% accuracy), which demonstrates the critical role of MC-based uncertainty estimation in stabilizing test-time model merging.
  \item \textbf{Insufficient perturbations}: For $p_{drop}=0.05$, the perturbations are too small to robustly separate stable and unstable parameter structures, yielding a sub-optimal accuracy of 90.42\%.
  \item \textbf{Optimal stability basin}: At $p_{drop}=0.10$, we obtain the optimal balance between gradient smoothing under expected prediction entropy and dynamic logit variance estimation, achieving our peak accuracy of 92.52\%.
  \item \textbf{Excessive masking noise}: When $p_{drop}$ is increased further to $0.20$ and $0.30$, accuracy decreases to 92.10\% and 91.15\% respectively. This drop is due to excessive masking of features, which injects severe noise into the self-supervised representation channel and degrades the representation quality under expected entropy optimization.
\end{enumerate}

\subsection{Robustness under Test-Time Covariate Shift}
To thoroughly evaluate the reliability and generalizability of MC-VTI, we perform a stress-test under simulated test-time covariate shifts. Specifically, we apply two distinct physical corruptions to the incoming images of the sequential test stream:
\begin{enumerate}
  \item \textbf{Gaussian Noise}: Additive white Gaussian noise with standard deviation $\sigma=0.4$ is added to each image, simulating severe sensor or transmission interference.
  \item \textbf{Dimming}: A constant brightness reduction of $-1.0$ is subtracted from the normalized pixel intensities, simulating severe under-exposure or dim lighting conditions.
\end{enumerate}

Table~\ref{tab:covariate_shift} summarizes the sequential stream accuracy (\%) of all compared methods under these out-of-distribution shifts.

\begin{table}[ht]
  \caption{Comparison of sequential test stream accuracy (\%) under simulated test-time covariate shifts.}
  \label{tab:covariate_shift}
  \vskip 0.15in
  \begin{center}
  \begin{small}
  \begin{sc}
  \begin{tabular}{lccc}
    \toprule
    Method & No Shift & Gaussian Noise ($\sigma=0.4$) & Dimming ($-1.0$) \\
    \midrule
    Static Merged  & 63.75 & 29.60 & 36.67 \\
    Standard TTA   & 93.56 & 59.10 & \textbf{81.00} \\
    S2C-Merge      & 92.56 & 41.67 & 77.85 \\
    EWC-TTA        & 81.19 & \textbf{60.40} & 74.10 \\
    \textbf{MC-VTI (Ours)} & \textbf{92.71} & 44.77 & 55.94 \\
    \bottomrule
  \end{tabular}
  \end{sc}
  \end{small}
  \end{center}
  \vskip -0.1in
\end{table}

We observe several highly compelling and honest scientific insights from these results:
\begin{enumerate}
  \item \textbf{Vulnerability to high-uncertainty drift}: Under severe covariate shifts, the input data becomes heavily degraded, which naturally spikes the test-time model uncertainty. In MC-VTI, this results in an extremely high stochastically estimated logit variance $V_c$. Consequently, the logit variance-weighted quadratic penalty applies a massive constraint on the classification heads, forcing them to stay close to their initial pre-trained parameters.
  \item \textbf{The double-edged sword of test-time dropout}: Crucially, under severe noise, evaluating the expected entropy across stochastically perturbed forward passes ($p_{drop}=0.1$) further exacerbates prediction flattening, injecting high-variance gradients into the self-supervised objective for updating the merging coefficients $\lambda$. This leads to poor weight-routing adjustments (sub-optimal $\lambda$), causing MC-VTI to perform worse than the frozen-classifier S2C-Merge baseline (44.77\% vs 41.67\% under noise, and 55.94\% vs 77.85\% under dimming), since S2C-Merge evaluates clean unperturbed images for updating $\lambda$.
  \item \textbf{Adaptation under unconstrained heads}: Conversely, unconstrained methods such as Standard TTA and EWC-TTA allow the classification heads to freely adjust their boundaries to adapt to the noisy or dimmed representations, recovering substantial accuracy (up to 60.40\% under noise and 81.00\% under dimming). However, this unconstrained adaptation is highly risky in standard sequential environments, as it is prone to catastrophic representation collapse if the stream has no shift.
\end{enumerate}

\end{document}
